\documentclass{article}

\usepackage{amsmath,amssymb}
\usepackage{xurl}
\usepackage[hidelinks]{hyperref}
\setlength{\emergencystretch}{2em}

% Shared commands for notes in docs/paper-gaps/.

\DeclareUnicodeCharacter{2080}{\ifmmode{}_{0}\else\textsubscript{0}\fi}
\DeclareUnicodeCharacter{2081}{\ifmmode{}_{1}\else\textsubscript{1}\fi}
\DeclareUnicodeCharacter{2082}{\ifmmode{}_{2}\else\textsubscript{2}\fi}
\DeclareUnicodeCharacter{2099}{\ifmmode{}_{n}\else\textsubscript{n}\fi}

\newcommand{\Lean}{\textsc{Lean}}
\ExplSyntaxOn
\NewDocumentCommand{\leanid}{m}{
  \ifmmode
    \textsf{\detokenize{#1}}
  \else
    \group_begin:
    \urlstyle{sf}
    \tl_set:Nn \l_tmpa_tl {#1}
    \tl_replace_all:Nnn \l_tmpa_tl {\_} {_}
    \exp_args:NV \path \l_tmpa_tl
    \group_end:
  \fi
}
\ExplSyntaxOff
\providecommand{\tr}{\operatorname{tr}}
\newcommand{\tnleanIssueBase}{https://github.com/LionSR/TNLean/issues/}
\newcommand{\tnleanPrBase}{https://github.com/LionSR/TNLean/pull/}
\newcommand{\tnleanDocBase}{https://sirui-lu.com/TNLean/docs/}
\newcommand{\ghissue}[1]{\href{\tnleanIssueBase#1}{GitHub issue \##1}}
\newcommand{\ghpr}[1]{\href{\tnleanPrBase#1}{PR \##1}}
\newcommand{\doclink}[1]{\href{\tnleanDocBase#1.html}{\path{#1}}}

% Dirac notation and the identity operator, as in blueprint/src/macros/common.tex.
% \providecommand keeps note-local definitions authoritative.
\usepackage{dsfont}
\providecommand{\ket}[1]{|#1\rangle}
\providecommand{\bra}[1]{\langle #1|}
\providecommand{\braket}[2]{\langle #1 | #2 \rangle}
\providecommand{\ketbra}[2]{|#1\rangle\!\langle #2|}
\providecommand{\Id}{\mathds{1}}
\providecommand{\one}{\mathds{1}}
\providecommand{\C}{\mathbb{C}}
\providecommand{\MN}[1]{M_{#1}(\C)}
\providecommand{\MD}{M_D(\C)}

% External referencing for paper sources.  The build helper writes sanitized
% auxiliary files under build/paper-aux/ and excludes duplicated source labels.
\usepackage{xr}

% Import a generated paper auxiliary file with a caller-supplied label prefix.
% The candidate paths support builds from docs/paper-gaps/, the repository root,
% and build/paper-aux/ itself.
\newcommand{\importpaperaux}[2]{%
  \IfFileExists{../../build/paper-aux/#2.aux}{%
    \externaldocument[#1]{../../build/paper-aux/#2}%
  }{%
    \IfFileExists{build/paper-aux/#2.aux}{%
      \externaldocument[#1]{build/paper-aux/#2}%
    }{%
      \IfFileExists{#2.aux}{%
        \externaldocument[#1]{#2}%
      }{}%
    }%
  }%
}

% General external reference macros
\newcommand{\paperref}[2]{\ref{#1-#2}}
\newcommand{\papereqref}[2]{(\ref{#1-#2})}

% CPSV16 source (1606.00608 / MPDO-22-12-17-2.tex) external document and shortcuts
\importpaperaux{cpsv16-}{1606.00608/MPDO-22-12-17-2-tnlean}
\newcommand{\cpsvref}[1]{\paperref{cpsv16}{#1}}
\newcommand{\cpsveqref}[1]{\papereqref{cpsv16}{#1}}

% Verdict marker: \gapnote{<kind>}{<status>}. Kind is the policy
% classification (clarification, local-correction, scope-restriction,
% unfaithful, false-source, open-gap); status is open, wip (elimination
% underway), resolved, or historical. Machine-read by the site index;
% prints a verdict line.
\newcommand{\gapnote}[2]{\par\noindent\textbf{Verdict:} #1 (#2).\par}


\title{Source-Notation Gap: The Majumdar--Ghosh Example Tensor}
\author{}
\date{2026-09-25}

\begin{document}
\maketitle
\gapnote{local-correction}{resolved}

\section{The assertion}

The review describes the Majumdar--Ghosh state as the ground state of the
spin-$\tfrac12$ chain
$H=\sum_i\mathbf S_i\cdot\mathbf S_{i+1}+\tfrac12\sum_i\mathbf S_i\cdot\mathbf S_{i+2}$,
the superposition of the two nearest-neighbour singlet coverings of a periodic
chain, and states that it ``can be written as an MPS with''
\cite[lines~2397--2405]{Cirac2021Matrix}
\begin{align}
    A
        &=
        \left[\ket{0}((02|+(20|)
            +
            \ket{1}((12|+(21|)\right]\otimes Y,
    \label{eq:rmp_majumdar_ghosh_tensor_gap_review_tensor}\\
    Y
        &=
        \begin{pmatrix}
            0 & -1\\
            1 & 0
        \end{pmatrix}.
    \label{eq:rmp_majumdar_ghosh_tensor_gap_singlet_matrix}
\end{align}
Here $Y$ is the singlet tensor of the AKLT example of the same appendix
\cite[lines~2380--2383]{Cirac2021Matrix}: $Y_{ab}$ is the coefficient of
$\ket{ab}$ in the singlet $\ket{10}-\ket{01}$. Read $(ab|$ as the bra of the pair
of virtual indices, so that $(ab|$ is the matrix unit $E_{ab}$ with row $a$ and
column $b$, and put
\begin{align}
    B^0
        &=
        E_{02}+E_{20},
    &
    B^1
        &=
        E_{12}+E_{21}.
    \label{eq:rmp_majumdar_ghosh_tensor_gap_auxiliary_matrices}
\end{align}

\section{The literal reading}

Read with $\otimes$ a Kronecker product, the formula
(\ref{eq:rmp_majumdar_ghosh_tensor_gap_review_tensor}) is the tensor
$\widetilde A^i=B^i\otimes Y$ on the six-dimensional bond
$\C^3\otimes\C^2$. Its periodic coefficients factor as
\begin{align}
    \tr\bigl(\widetilde A^{i_1}\cdots\widetilde A^{i_N}\bigr)
        &=
        \tr\bigl(B^{i_1}\cdots B^{i_N}\bigr)\,\tr\bigl(Y^N\bigr),
    \label{eq:rmp_majumdar_ghosh_tensor_gap_literal_factorization}
\end{align}
and $Y^2=-I_2$. The level $2$ of the three-level bond is a hub: no matrix
$B^i$ connects the levels $0,1$ with each other, and none maps $2$ to itself.
A closed path of length $N=2m$ therefore alternates between the hub and the
other levels, and a pair of consecutive letters $a,b$ returning to the hub has
amplitude $(B^aB^b)_{22}=\delta_{ab}$. Writing $\Phi_{\mathrm{even}}$ and
$\Phi_{\mathrm{odd}}$ for the coverings $(1,2)(3,4)\cdots(N-1,N)$ and
$(2,3)(4,5)\cdots(N,1)$ by the pair state $\ket{00}+\ket{11}$,
\begin{align}
    \ket{V^{(2m)}(\widetilde A)}
        &=
        2(-1)^m\bigl(\ket{\Phi_{\mathrm{even}}}+\ket{\Phi_{\mathrm{odd}}}\bigr).
    \label{eq:rmp_majumdar_ghosh_tensor_gap_literal_vector}
\end{align}
On the configuration with all spins $0$ this vector has coefficient
$4(-1)^m\neq0$, whereas every singlet covering vanishes there. The literal
tensor therefore does not represent the Majumdar--Ghosh state at any even
length: its periodic vector is not a scalar multiple of the dimer
superposition.

A gauge transformation of the bond leaves the periodic vector unchanged, so the
only freedom left is a change of basis of the physical spin, a single-site
matrix $U$ applied at every site. It does not help either. The vector
$U^{\otimes N}(\ket{\Phi_{\mathrm{even}}}+\ket{\Phi_{\mathrm{odd}}})$ is the sum
of the two coverings by the pair state
$(U\otimes U)(\ket{00}+\ket{11})=\sum_{ab}M_{ab}\ket{ab}$ with the symmetric
coefficient matrix $M=UU^{\mathsf T}$. Suppose it equals $c$ times the sum of the
two coverings by the singlet $\ket{10}-\ket{01}$, whose coefficient matrix is
$Y$. Compare coefficients on four configurations of the ring of $N=2m$ sites.
On the configuration with all spins $0$ the left side is $2M_{00}^m$ and the
right side vanishes, so $M_{00}=0$; likewise $M_{11}=0$. Put
$\mu=M_{01}=M_{10}$. On the alternating configuration $0101\cdots01$ both
coverings pair a $0$ with a $1$, and
\begin{align}
    2\mu^m
        &=
        c\bigl(Y_{01}^m+Y_{10}^m\bigr)
        =
        c\bigl((-1)^m+1\bigr).
    \label{eq:rmp_majumdar_ghosh_tensor_gap_alternating}
\end{align}
If $m$ is odd, then $\mu=0$. If $m\geq2$ is even, then $c=\mu^m$. Now take
the configuration $0110$ followed by $m-2$ copies of $01$. The covering
$(1,2)(3,4)\cdots$ pairs $01$, $10$, and then $01$ repeatedly, while the
covering $(2,3)(4,5)\cdots$ contains the pair $11$. Hence
$\mu^m=cY_{01}Y_{10}Y_{01}^{m-2}=-c$, and together with $c=\mu^m$ this gives
$\mu=0$. In either case $M=0$, so $U$ is singular and the left side
vanishes, which forces $c=0$. No change of physical basis turns the literal
tensor into a tensor for the Majumdar--Ghosh state: the symmetric pair state is
the obstruction.\footnote{Formal declarations: \leanid{MPSTensor.majumdarGhoshReviewTensor},
\leanid{MPSTensor.majumdarGhoshReview\_coeff},
\leanid{MPSTensor.majumdarGhoshReview\_mpv\_eq}, and
\leanid{MPSTensor.majumdarGhoshReview\_mpv\_ne\_smul} in
\doclink{TNLean/MPS/Examples/MajumdarGhoshDimer}.}

\section{The correction: the singlet on the spin levels}

The intended construction places the singlet $Y$ on the two spin levels $0,1$
of the three-level bond and carries the empty level $2$ through, that is, it
follows each $B^i$ by the bond state $Y\oplus1$:
\begin{align}
    C^i
        &=
        B^i\,(Y\oplus1),
    &
    C^0
        &=
        E_{02}-E_{21},
    &
    C^1
        &=
        E_{12}+E_{20}.
    \label{eq:rmp_majumdar_ghosh_tensor_gap_corrected_tensor}
\end{align}
The level $2$ is again a hub, and now the two-letter hub amplitude is
$(C^aC^b)_{22}=Y_{ab}$. Hence, for every even $N>0$,
\begin{align}
    \ket{V^{(N)}(C)}
        &=
        \ket{S_{\mathrm{even}}}+\ket{S_{\mathrm{odd}}},
    \label{eq:rmp_majumdar_ghosh_tensor_gap_corrected_vector}
\end{align}
where $S_{\mathrm{even}}$ and $S_{\mathrm{odd}}$ are the two coverings by the
singlet $\ket{10}-\ket{01}$. This is the review's superposition. The
correction replaces the Kronecker factor $Y$ in
(\ref{eq:rmp_majumdar_ghosh_tensor_gap_review_tensor}) by the bond state
$Y\oplus1$ applied after $B^i$; the bond dimension is three, not
six.\footnote{Formal declarations:
\leanid{MPSTensor.majumdarGhoshBondSingletTensor} and
\leanid{MPSTensor.majumdarGhoshBondSinglet\_mpv\_eq\_pairCovering} in
\doclink{TNLean/MPS/Examples/MajumdarGhoshDimer}.}

\section{The formal representative}

The formal example uses the bond-dimension-three tensor
\begin{align}
    A^0
        &=
        \begin{pmatrix}
            0 & 1 & 0\\
            0 & 0 & 0\\
            1/\sqrt{2} & 0 & 0
        \end{pmatrix},
    &
    A^1
        &=
        \begin{pmatrix}
            0 & 0 & 1\\
            -1/\sqrt{2} & 0 & 0\\
            0 & 0 & 0
        \end{pmatrix},
    \label{eq:rmp_majumdar_ghosh_tensor_gap_formal_matrices}
\end{align}
whose hub is the level $0$, with two-letter hub amplitude
$(A^aA^b)_{00}=Y_{ab}/\sqrt2$, the normalized singlet. For every even $N>0$
its periodic vector is the sum of the two coverings by normalized singlets,
and $\ket{V^{(N)}(C)}=\sqrt2^{\,N/2}\ket{V^{(N)}(A)}$. On odd rings the
periodic vector vanishes.

Each covering, not only their sum, lies in the periodic three-site chain
ground space of $A$ for even $N\ge4$: with $G=\operatorname{diag}(1,-1,-1)$,
which anticommutes with $A^0$ and $A^1$, the vectors
$\sigma\mapsto\tr(A^{\sigma_1}\cdots A^{\sigma_N})$ and
$\sigma\mapsto\tr(A^{\sigma_1}\cdots A^{\sigma_N}G)$ both satisfy every window
constraint, and their half-sum and half-difference are the two coverings. The
three-site ground space of $A$ equals the eigenspace of
$h=\tfrac12(\mathbf S_1\cdot\mathbf S_2+\mathbf S_2\cdot\mathbf S_3
+\mathbf S_1\cdot\mathbf S_3)$ for the eigenvalue $-\tfrac38$, the three-spin
states without a total-spin-$\tfrac32$ component. Consequently every vector of
the chain ground space, and in particular each covering, satisfies
$H\psi=-\tfrac{3N}{8}\psi$.\footnote{Formal declarations:
\leanid{MPSTensor.majumdarGhosh\_mpv\_eq\_pairCovering},
\leanid{MPSTensor.majumdarGhoshBondSinglet\_mpv\_eq\_smul},
\leanid{MPSTensor.majumdarGhosh\_groundSpace\_three\_eq\_eigenspace},
\leanid{MPSTensor.majumdarGhosh\_pairCoveringEven\_mem\_chainGroundSpace},
\leanid{MPSTensor.majumdarGhosh\_pairCoveringOdd\_mem\_chainGroundSpace}, and
\leanid{MPSTensor.majumdarGhoshHamiltonian\_apply\_of\_mem\_chainGroundSpace} in
\doclink{TNLean/MPS/Examples/MajumdarGhoshDimer} and
\doclink{TNLean/MPS/Examples/MajumdarGhoshHamiltonian}.}

\section{Remaining formal boundary}

Two statements are not formalized. The first is the lower bound
$H\ge-\tfrac{3N}{8}$, which follows from
$H=\sum_i h_i$ and $h\ge-\tfrac38$ and makes the eigenvectors above ground
states. The second is that the periodic three-site chain ground space of $A$
is exactly two-dimensional on even rings of length at least four, spanned by
the two coverings. Neither affects the notational correction recorded here.

\section{Verdict}

This is a local correction of the source notation. Read literally, the
review's formula gives a six-dimensional tensor whose periodic vector is the
superposition of two coverings by $\ket{00}+\ket{11}$ rather than by singlets.
With the singlet $Y$ placed on the spin levels of a three-level bond, the
formula gives exactly the Majumdar--Ghosh superposition, and the periodic
vectors of the formal bond-dimension-three representative are those of the
corrected tensor multiplied by $\sqrt2^{\,-N/2}$.

\bibliographystyle{alpha}
\bibliography{references}

\end{document}
